% !TEX program = xelatex
% ============================================================
%  main_zh.tex — 第四章「方法」中文稿（与 main.tex 对应）
%  图片与 main.tex 相同，放在本文件同级目录
%
%  编译器（二选一，按你日志「仍在用 pdfLaTeX」说明）:
%  A) 推荐：主菜单/设置里把 Compiler 改为 XeLaTeX（与上一行 magic comment 一致）
%  B) 若只能 pdfLaTeX：本稿已用 fontset=ubuntu（文鼎，避免 fandol 报 unavailable
%     后回退到 unisong9e 崩掉）。需完整 TeX Live 中文包，如仍缺字执行:
%        tlmgr install arphic-ttf cjk
%     或安装发行版的 texlive-lang-chinese / collection-langcjk
% ============================================================
\documentclass[UTF8,11pt,a4paper,fontset=ubuntu]{ctexart}

\usepackage[margin=1in]{geometry}
\usepackage{graphicx}
\usepackage{amsmath,amssymb,amsthm,bm}
\usepackage{booktabs}
\usepackage{multirow}
\usepackage{array}
\usepackage{siunitx}
\usepackage[hidelinks]{hyperref}
\usepackage{cite}
\usepackage{caption}
\usepackage{subcaption}
\usepackage{algorithm}
\usepackage{algpseudocode}
\usepackage{xcolor}
\usepackage{microtype}

\graphicspath{{./}}

\newcommand{\Pof}{\mathbf{P}^{\mathrm{OOF}}}
\newcommand{\logit}{\operatorname{logit}}
\newcommand{\RR}{\mathbb{R}}
\newcommand{\EE}{\mathbb{E}}
\newcommand{\ptilde}{\tilde p}
\newcommand{\phat}{\hat p}
\newcommand{\taustar}{\tau^{\star}}
\DeclareMathOperator*{\argmin}{arg\,min}
\DeclareMathOperator*{\argmax}{arg\,max}

\title{可执行的员工离职预测系统：\\
面向应届生的噪声鲁棒堆叠与硬约束反事实干预}
\author{Yuhang Song \and 合作者 \\ \small{\today}}
\date{}

\begin{document}
\maketitle

\section{方法}
\label{sec:methods}

\subsection{概述}
\label{sec:methods-overview}

我们将整个系统视为两个紧耦合的组成部分：
（i）一条\textbf{推理流程}，将单条员工原始记录映射为二值风险预警；
（ii）一套\textbf{训练协议}，估计推理阶段所需的全部参数。
图~\ref{fig:inference} 展示前者，图~\ref{fig:training} 展示后者。两图由同一特征映射
$\varphi(\cdot)$、同一样本权向量 $\mathbf{w}$，以及同一组冻结的导出物
$\{\text{base}_k,\bm{\beta},b,g,\taustar\}$ 相联系；该集合中每个对象在推理图中
有唯一使用位置，在训练流程中也有唯一产生位置。

\subsection{特征变换 \texorpdfstring{$\varphi$}{phi}}
\label{sec:methods-phi}

设 $\mathbf{x}\in\RR^{12}$ 为十二条整数编码调研变量的原始输入（三列人口学特征、
四列工作特征、五列工作情境特征）。我们施加确定性变换
$\varphi:\RR^{12}\to\RR^{41}$，将六类互补信息拼接为一条特征向量：
（1）在有序题项上构造十二条累积型 Likert 哑变量
$\mathbf{1}\{x_j\ge c\}$，$c\in\{2,3,4,5\}$；
（2）十二条原始数值特征；
（3）六条人工构造的二元交互（例如收入 $\times$ 压力、满意度 $\times$ 匹配等）；
（4）八条「期望--现实」落差特征；
（5）对收入作 $\log(\text{income})$ 与 $\text{income}^2$ 两项变换；
（6）单条经 K 折目标编码的复合分数。
树模型直接使用 $\varphi(\mathbf{x})$；基于距离与梯度的学习器在 $\varphi(\mathbf{x})$
上额外做 $z$ 分标准化
$\tilde{\varphi}(\mathbf{x}) = (\varphi(\mathbf{x}) - \bm{\mu}_{\mathrm{train}})/
\bm{\sigma}_{\mathrm{train}}$，
其中 $\bm{\mu}_{\mathrm{train}},\bm{\sigma}_{\mathrm{train}}$ 仅在有标签的训练折上估计，以避免信息泄漏。

\subsection{推理流程}
\label{sec:methods-inference}

推理流程（图~\ref{fig:inference}）在概念上可分成六个阶段。
给定输入 $\mathbf{x}$，五个基学习器
$k\in\{\mathrm{RF}, \mathrm{NRB}, \mathrm{MT}, \mathrm{SVM}, \mathrm{ET}\}$
分别输出概率 $p_k = f_k(\varphi(\mathbf{x}))\in[0,1]$。
一级（Level-1）L2 正则化逻辑回归元学习器在 logit 空间上融合五个基模型输出：
\begin{equation}
\label{eq:meta}
\ptilde \;=\; \sigma\!\left(\sum_{k=1}^{5}\beta_k\,\logit(p_k) + b\right),
\qquad
\sigma(z)=\frac{1}{1+e^{-z}}.
\end{equation}
将元模型输出 $\ptilde$ 经\textbf{后验}等距校准（isotonic）映射
$g:[0,1]\to[0,1]$（用 Pool-Adjacent-Violators 算法学习得到的单调非减阶梯函数），
得到已校准风险分数
\[
\phat \;=\; g(\ptilde).
\]
再使用在验证准则下选取的平衡准确率最优阈值 $\taustar$ 将 $\phat$ 二值化：
\[
\hat y \;=\; \mathbf{1}\{\phat \ge \taustar\}, \qquad \taustar = 0.135.
\]
上述右端量均在部署时固定；\ref{sec:methods-training}~节说明其估计方式。

\begin{figure}[t]
  \centering
  \includegraphics[width=\linewidth]{fig_A_inference.pdf}
  \caption{\textbf{所提出堆叠集成的推理流程。}
  原始输入 $\mathbf{x}\in\RR^{12}$ 经确定性映射
  $\varphi:\RR^{12}\!\to\!\RR^{41}$ 后分为两支：树模型（RF、NR-Boost、ET）直接使用
  $\varphi(\mathbf{x})$，距离/梯度类学习器（SVM-RBF、MT-MLP）使用其 $z$ 分标准化
  版本。每个基学习器 $k$ 输出概率 $p_k\in[0,1]$（彩色圆点）。L2 正则逻辑回归
  元学习器在 logit 空间上融合五路输出（式~\ref{eq:meta}）；边颜色表示 $\beta_k$ 的符号
 （蓝为正、红为负），边宽表示 $|\beta_k|$。$\ptilde$ 经 isotonic 映射 $g$ 得到
  已校准风险 $\phat$，再由平衡准确率意义下的最优阈值 $\taustar=0.135$ 得到二值预警
  $\hat y$。各箭头旁标注张量形状；圈注数值示例展示一名高风险员工的完整前向计算轨迹。}
  \label{fig:inference}
\end{figure}

\subsection{训练协议}
\label{sec:methods-training}

训练协议（图~\ref{fig:training}）说明各部署产物的估计方式，分四个阶段。

\paragraph{阶段 1：Cleanlab-v6 样本权。}
为缓解自陈式离职调研中普遍存在的标签噪声，我们在独立交叉验证回路上拟合
噪声鲁棒探针分类器，为每个训练样本计算自置信（self-confidence）分数；按分数升序
将\textbf{末位 10\%} 标为可能错标，赋降权 $w_i=0.3$，其余 90\% 保持 $w_i=1.0$。
得到二元权向量 $\mathbf{w}\in\RR^{N}$（$N=4{,}375$）并落盘，供后续各步复用。

\paragraph{阶段 2：5 折 OOF 训练。}
对每个基学习器 $f_k$，在层化 5 折交叉验证中通过原生 \texttt{sample\_weight}
接口注入 $\mathbf{w}$；每折在四个入折上训练、在一个出折上预测，将五折出折
预测拼接，得到无泄漏的 OOF 概率矩阵
\[
\Pof \;=\; \bigl[\,p_k^{\mathrm{OOF}}(\mathbf{x}_i)\,\bigr]_{i,k}
\;\in\; [0,1]^{N\times 5}.
\]
该矩阵即元学习器的训练输入；由构造，$p_k^{\mathrm{OOF}}(\mathbf{x}_i)$ 所对应的
基模型在拟合时未见过样本 $i$，从而避免朴素堆叠中常见的一级（Level-1）泄漏。

\paragraph{阶段 3：元模型拟合与 isotonic 校准。}
对 $\Pof$ 的每一列作 logit 变换，得到元特征
$\mathbf{l}_i = \logit(\mathbf{p}_i^{\mathrm{OOF}})\in\RR^{5}$，
并最小化带 L2 惩罚的加权交叉熵以拟合 L2 正则逻辑回归：
\begin{equation}
\label{eq:meta-objective}
\{\hat{\bm\beta},\hat b\}
\;=\;
\argmin_{\bm\beta,\,b}\;
-\sum_{i=1}^{N} w_i
\bigl[\,y_i\ln\hat p_i + (1-y_i)\ln(1-\hat p_i)\,\bigr]
\;+\;\frac{1}{2C}\|\bm\beta\|_2^{2},
\end{equation}
其中 $\hat p_i = \sigma(\bm\beta^{\!\top}\mathbf{l}_i + b)$，$C=10$。目标为凸，可用
LBFGS 求解。随后在各基学习器上\textbf{以全量训练集}（带 $\mathbf{w}$）重拟合
得到可部署模型，并计算其 OOF 模式下的元输出 $\ptilde_{\mathrm{oof},i}$。在
$(\ptilde_{\mathrm{oof},i}, y_i)$ 上再拟合 isotonic 校准 $g$（Pool-Adjacent-Violators）以保证单调。决策阈值取为
\[
\taustar \;=\; \argmax_{\tau\in(0,1)}\;
\mathrm{BalAcc}\!\bigl(\{\phat_{\mathrm{oof},i}\}, y\,;\tau\bigr),
\qquad \taustar = 0.135,
\]
其中 $\phat_{\mathrm{oof},i} = g(\ptilde_{\mathrm{oof},i})$。此后 $g$、$\taustar$ 与
$(\hat{\bm\beta},\hat b)$ 均\textbf{冻结}，不再在测试集上重估。

\paragraph{阶段 4：部署产物包。}
最终导出：（a）五个重拟合后的基学习器 $\{f_k\}_{k=1}^{5}$；（b）元系数
$(\hat{\bm\beta},\hat b)$；（c）isotonic 映射 $g$；（d）阈值 $\taustar$。它们恰为
图~\ref{fig:inference} 推理各节点所消费的四类对象。

\begin{figure}[t]
  \centering
  \includegraphics[width=\linewidth]{fig_B_training.pdf}
  \caption{\textbf{冠军模型的训练协议。}
  \textbf{阶段 1（Cleanlab 权）。} 噪声鲁棒探针在独立 CV 上计算逐样本自置信；末位
  10\% 设 $w_i=0.3$，其余 $w_i=1.0$，得到 $\mathbf{w}$。
  \textbf{阶段 2（5 折 OOF）。} 五个基学习器在带 $\mathbf{w}$ 的层化 5 折中训练，拼接
  出折预测得到无泄漏的 OOF 矩阵 $\Pof\in[0,1]^{N\times 5}$。
  \textbf{阶段 3（元模型 + Isotonic）。} 对 $\Pof$ 作 logit 后送入 L2 正则逻辑回归，
  最小化式~\ref{eq:meta-objective}（黄框，LBFGS）得到 $(\hat{\bm\beta},\hat b)$；各
  基模型再在全量训练集上重拟合；在 $(\ptilde_{\mathrm{oof}}, y)$ 上拟合 isotonic
  $g$；在 OOF 预测上按平衡准确率选取 $\taustar$。
  \textbf{阶段 4（冻结产物）。} 导出物 $\{f_k,\hat{\bm\beta},\hat b,g,\taustar\}$ 与
  图~\ref{fig:inference} 推理时所用完全一致。}
  \label{fig:training}
\end{figure}

\subsection{与图~\ref{fig:inference}、图~\ref{fig:training} 的对应关系}
\label{sec:methods-bridge}

图~\ref{fig:inference} 与图~\ref{fig:training} 是同一套模型的两个互补视角：
前者在参数已固定的前提下，展示单条输入如何经冻结网络得到二值预警；后者则揭示
可学习量——五类基学习器内部参数、元系数 $(\bm\beta,b)$、isotonic 映射 $g$ 与阈值
$\taustar$——如何从数据中间接估计。两图共享（i）同一 41 维特征空间 $\varphi$；
（ii）同一样本权 $\mathbf{w}$；（iii）同一组基学习器超参数；（iv）同一 5 折划分。
图~\ref{fig:training} 阶段~4 产出的每个对象在图~\ref{fig:inference} 中有且仅有一个
使用位置：$(\bm\beta,b)$ 仅出现在 meta-$\sigma$ 节点，$g$ 仅出现在 isotonic 节点，
$\taustar$ 仅出现在阈值节点。这种训练--推理一一对应使复现可归结为单一产物包
完整性校验，且阶段~2 的 OOF 构造保证元学习器在训练时见到的基模型概率分布
与部署时一致。

\subsection{硬约束反事实（HRCF）干预}
\label{sec:methods-hrcf}
% TODO — 见 docs/paper/06_methods_hrcf.md
\textit{（待写）在强可行动约束下基于投影的反事实生成；代价函数；top-$k$ 方案枚举；
与软约束型 DiCE 变体的比较。}

\end{document}
